\section{The Condor and the Anarch}

\begin{quotex}
We fail not because of our dreams but because we do not dream forcefully enough. \flright{\textsc{Ernst Junger}}

The Rose does not ask why. It blooms because it blooms. \flright{\textsc{Angelus Silesius}}

I live after a calling as little as the flower grows and gives fragrance after a calling. \flright{\textsc{Max Stirner}}

Of course, the anarch is hemmed in from the first day by father and mother, by state and society. Those are prunings, tappings of the primordial strength, and nobody escapes them. One has to resign oneself. But the anarchic remains, at the very bottom, as a mystery, usually unknown even to its bearer. It can erupt from him as lava, can destroy him, liberate him. Distinctions must be made here: love is anarchic, marriage is not. The warrior is anarchic, the soldier is not. Manslaughter is anarchic, murder is not. Christ is anarchic, Saint Paul is not. Since, of course, the anarchic is normal, it is also present in Saint Paul, and sometimes it erupts mightily from him. Those are not antitheses but degrees. The history of the world is moved by anarchy. In sum: the free human being is anarchic, the anarchist is not. \flright{\textsc{Ernst Junger}, \emph{Eumeswil}}

\end{quotex}
While I was preparing a book review of \emph{Eumeswil} by \textbf{Ernst Junger}, Keith Preston published this excellent review of a new translation of Max Stirner's \textit{Der Einzige und sien Eigentum}\footnote{\url{https://attackthesystem.com/2017/07/24/a-little-less-piousness-please/}}. I use the older translation titled in English \emph{The Ego and his Own}, but the new one is titled \emph{The Unique and its Property}. Apparently there is no concept in English that captures the meaning of the German.

Since Eumeswil is more or less an extended meditation on Stirner's book, I had to go back and make some revisions. Stirner is of interest to us, not least because of his influence on Carl Schmitt, Julius Evola, and even Rudolf Steiner who penned the \emph{Philosophy of Freedom} as the metaphysical justification for Stirner's views. Steiner once described himself as an Individualist Anarchist\footnote{\url{http://philosophyoffreedom.com/profiles/blog/list/tag/max+stirner}}. Evola wrote a book on Junger's idea of the Worker and his notion of the Autarch seems close to the latter's Anarch. The Anarch also rides the tiger in a world he can neither consent to nor oppose.

Eumeswil is a small post-cataclysm city-state located in the Maghreb, at an unspecified future date. The Condor is the nickname for its tyrant (as opposed to a despot). The narrator, Martin Venator, is a courtier in the Condor's Casbah by night, but an historian and scholar by day. This is not a conventional novel, but more like a series of meditations akin to Rousseau's \emph{Reveries of a Solitary Walker}. Presumably Venator's voice is also Junger's.

The name Martin comes from Mars, the god of war, and Venator means ``hunter" in Latin. Hunting is a recurring theme in the novel. ``Martin" is the name given by Venator's father, a point emphasized, and apparently not harmonious with his last name. Mars Hunter is rather aggressive, and Venator does not accept his natural father's view of life. However, the Condor renamed him ``Manuel", which is Hebrew for ``God with us". So henceforth, the Hunter is associated with a different God.

Venator had two primary mentors, one an historian, the other a philosopher. Vigo is modelled after \textbf{Giambattista Vico} whose guiding principle was ``The truth is the fact itself". Venator is interested solely in facts, not opinions or metaphysical speculations. On the other hand, Vico's views are balanced by Bruno, modeled after \textbf{Giordano Bruno}. A third factor is the grammarian who lectures on the importance of the use of words to legal issues. At a time like ours, grammar is being deliberately mangled in order to alter legal reasoning.

\paragraph{The Anarch}
Although Max Stirner is often associated with Left-wing anarchism, Venator rejects anarchism in favour of the Anarch. Hence, the book is an extended meditation in the difference between the two outlooks. The anarchist rejects authority and attacks the system, and the Anarch works within and around it, while being detached from it. The former tries to be outwardly free, the latter is inwardly free. Venator writes:

\begin{quotex}
It is not that I as an anarch reject authority à tout prix. On the contrary, I seek it, and that is precisely why I reserve the rights to examine it. 

\end{quotex}
There is a reason for this:

\begin{quotex}
The anarchist, as the born foe of authority, will be destroyed by it after damaging it more or less. The anarch, on the other hand, has appropriated authority; he is sovereign. He therefore behaves as a neutral power vis-à-vis state and society. He may like, dislike, or be indifferent to whatever occurs in them. That is what determines his conduct; he invests no emotional values. 

\end{quotex}
The anarchist scorns the rules, ``like people who deliberately drive on the wrong side of the road."

\begin{quotex}
The anarch, in contrast, knows the rules. He has studied them as a historian and goes along with them as a contemporary. Whenever possible, he plays his own game within their framework; this makes the fewest waves. 

\end{quotex}
\paragraph{The Anarchist}
The anarchist generally holds to these principles:

\begin{itemize}
\item \textbf{Atheism}: There is no necessary being, therefore everything is contingent. In practice, this means that there is no sufficient reason for any given state of affairs. Hence, any other possible state of affairs has its own legitimacy. 
\item \textbf{Antilogos}: A corollary is that there is no overarching order to the world. An imposed order, even by a god, is an affront to the freedom of the anarchist. 
\item \textbf{Positivism}: Reality can only be known through the senses. There is no transcendent source of knowledge. Any such claim is merely a human construct, or, in Stirner's terms, and airy ``spook". 
\item \textbf{Nominalism}: There is no abstract realm of pure ideas. Words are arbitrary human constructions. This goes by the name today as ``anti-essentialism". For example, there is no essence to being a man or a woman. Such characteristics are for the Ego or Unique One to determine. 
\item \textbf{Egoism}: The purpose of mind is to pursue the alleged self-interest of the being. Not surprisingly, this almost always reduces to sex, as in \textit{Kidz, Kulture, and Kreativity}\footnote{\url{https://gornahoor.net/?p=1582}}. Stirner and his subsequent acolytes felt, and still feel, obliged to flout bourgeois sexual morality, like people who deliberately drive on the wrong side of the road. When combined with anti-essentialism, then even ``biological facts" are considered an affront to freedom. 
\end{itemize}
Traditional ideas like ``tradition, nation, race, family, religion, morality, duty, order, obedience" refer to nothing in reality. Rather, they are ``spooks", i.e., human creations, a superstructure added to the natural world. The scientific view is that they are mere bigoted opinions\footnote{\url{https://gornahoor.net/?page_id=8496}}.

The anti-philosophy of nominalism, in practice, means that the most powerful and most wealthy are able to control the discourse, or ``spooks" if you prefer. There is no concern for objective truth, just for the satisfaction of desires and will. Hence, if science is objective, nominalism is not scientific. Only realism is scientific and thus able to effectively combat the powerful. The man who understands reality does not believe in spooks, i.e., he knows what is real and what is an artifact.

\paragraph{The Luminar}
Venator had access to the Luminar, a strange brew of Wikipedia, the Akashic Record, and spirit channeling. Through it, he could see the unfolding of the whole scope of human history. Hence, he couldn't consent to any particular movement, since each one would pass. Only facts counted. What is considered right or legal can change.

Apparently, the luminar could reconstruct historical personalities so that Venator could engage in discussions with Bachofen on matriarchy and with Moses on morals and religion. Although Venator is free from ``moral and religious bonds", there are some problems that he ``can scarcely resolve, but that have an effect precisely by surfacing" — for example, the Resurrection.

\paragraph{Three Great Saints}
\begin{quotex}
One consequence of the worldwide entanglement is that solitary men appear, talents not rooted in one specific landscape or tradition. 

\end{quotex}
Bruno taught Venator detachment. He points out that the view that external things like ``rank, money, and honors" bring happiness, is not necessarily correct. Thomas Aquinas would describe them as ``accidental" qualities. If so, then what exactly is essential to the Unique One? What is his Own or his Property? The Ego has a calling higher than the pursuit of money, sex, and power. Venator explains:

\begin{quotex}
If one manages to separate essence from flesh, if one manages, that is, to get distance from oneself, then one climbs the first step toward spiritual power. Many exercises are geared to this—from the soldier's drill to the hermit's meditation. 

\end{quotex}
Why does this elude the great majority of human beings? It is because they have no depth:

\begin{quotex}
Why do they not know that the world remains inalterable in change? Because they never find their way down to its real depth, their own. That is the sole place of essence, safety. 

\end{quotex}
So who shows us that depth?

\begin{quotex}
The characteristic feature of the great saints is that they get at the very heart of the matter. The most obvious things are invisible because they are concealed in human beings; nothing is harder to evince than what is self-evident. Once it is uncovered or rediscovered, it develops explosive strength. \textbf{Saint Anthony} recognized the power of the solitary man. \textbf{Saint Francis} that of the poor man. \textbf{Stirner} that of the ``only" man. At bottom, everyone is solitary, poor, and ``only" in the world. 

\end{quotex}
It is certainly curious to include Stirner among the great saints, ``of whom there are very few". What exactly is the matter whose heart Stirner reached? Venator mentions two of Stirner's axioms:

\begin{enumerate}
\item That is not My business. 
\item Nothing is more important than I 
\end{enumerate}
Stirner made the self-evident, evident, namely, in Venator's words:

\begin{quotex}
Each person is the center of the world, and his unconditional freedom creates the gap in which respect and self-respect balance out. 

\end{quotex}
For Stirner, there is nothing to achieve. He writes:

\begin{quotex}
As this rose is a true rose to begin with, this nightingale always a true nightingale, so I am not for the first time a true man when I fulfill my calling, live up to my destiny, but I am a ``true man" from the start. My first babble is the token of the life of a ``true man", the struggles of my life are the outpourings of his force, my last breath is the last exhalation of the force of the ``man". The true man does not lie in the future, an object of longing, but lies, existent and real, in the present. 

\end{quotex}
A dog does not aspire to be a ``good" dog, it has no essence as a task. Stirner applies this to man, also. For a man to aspire to be good, or rational, etc., is to aspire to be other than what he is. For Stirner, the fact of man is the exercise of force or strength. So who has the most force? Obviously, the man who is most conscious, who has self-mastery. Can I propose that as a goal? Not in Stirner's eyes, but some men will choose it, just as the rose chooses to bloom. It is an error to presume that all men are alike.

Venator seems, however, to have seen even more deeply, viz., the distinction between the anarch and the anarchist. He takes the first term, rejecting anarchy, while Stirner's followers mostly embrace it. Venator, the historian, needed to determine where anarchy's self-understanding in acting, thinking, or poetic creation occurred. Specifically, he sought the point where it coincided with the attainment of self-comprehension, the basis of freedom. Venator looked to the pre-Socratics (as did Heidegger), the Gnostics (as did Jung), and the mysticism of \textbf{Angelus Silesius}. Silesius was a German Catholic priest who continued the Nordic tradition that began with Meister Eckhart through Jacob Boehme, including John of Ruysbroek among others. Like the rose that blooms, the Anarch prays, the anarchist does not. There is no purpose to the prayer, it is simply his nature.

So we see that the basis of freedom begins in self-knowledge (or self-comprehension). After all, if nothing is more important than ``I", then that ``I" must be known. Even God commands you to love yourself, as if you needed that nudge. The traditional order of charity is: God, Self, and then spouse, children, parents, siblings, friends, and so on. Even the atheist can agree to most of that.

Venator seeks to transcend the material, perhaps because matter limits his freedom. One day he came to this realization:

\begin{quotex}
I thus managed to achieve something I had always dreamed of: a complete detachment from my physical existence. I saw myself in the mirror as a transcendent suitor—and myself, confronting him, as his fleeting mirror image. 

\end{quotex}
Angelus Silesius also found the transcendent suitor:

\begin{quotex}
I am like God and God like me.

I am as Large as God, He is as small as I.

He cannot be above me, nor I beneath him. 

\end{quotex}
Venator says that the anarch is to the anarchist as the monarch is to the monarchist. That is, the monarchist wishes there were a king, the monarch really is the king. Analogously, the anarchist wishes he were free, the anarch actually is free.

\paragraph{Postscript on the Condor}
Venator posed as a bird-watcher in order to have an excuse to wander in the outskirts of the city. He tells this tale of the Condor. A South American tribe would capture a condor and starve it for a week. Then they would tie the condor to the back of a bull that had been severely bloodied through deliberate cuts. The hungry condor would then rip out the flesh of the still breathing bull.

For more on spooks from 2010, see \textit{The Demon of Dialectics}\footnote{\url{https://gornahoor.net/?p=813}}.

\flrightit{Posted on 2017-07-31 by Cologero }

\begin{center}* * *\end{center}

\begin{footnotesize}\begin{sffamily}

\texttt{Max on 2017-08-05 at 16:18 said: }

Please allow me to re-germanize the English in order to produce an accurate translation of the title that also manages to preserve the artful wordplay intact: ``The Only and His Alonedom" [The One-Like and His All-One-Dom]. 

Dom is an archaic English word that has fallen out of use except in relic composite words such as kingdom, wisdom and freedom, and is the exact translation of the German suffix -tum, meaning judgement. No latin neologism to be seen here. As it happens, neologism is not possible within purely germanic tongue because of its fluid poetical nature, but is rather a phenomenon resulting from the interplay of the latin and the germanic spirit. (English is germanic spoken in latin mode, French is latin spoken in germanic mode. Thus, learning Latin, French, English and German should provide the whole spectral loom to understand the particular and recurrent modifications of the Roman current up until modern times.) Wisdom as an example means wise decree, so it is not something you ``have" in an abstract way, but an exercised activity. Likewise, property in the sense of eigentum or alonedom is actively exercised and enforced by an only ``alone" as implicated in the word.

In germanic languages the word for ``property" (in the sense of belongings) has the same linguistic form as kingdom. Perhaps that linguistic privation of modern English says something of the Anglo's and a presumptive jacobinism. I do not usually work myself up emotionally over trivial discussions, but recently got into quite a heated argument with an anarchist precisely over the nature of ``eigentum" when the inner anarch was unexpectedly provoked by the repulsive apparition of anarchism. 

It is curious that in popular English usage a positive word for exerting judgement has acquired associations such as ``doom and gloom". I guess that people are generally aversive of judgements, distinctions, discriminations and decisions, which by the linguistic evidence is an already hundreds of years old development. 

Now in my deeming, doom, in all of its wondrous forms – some of which listed above – cannot really come too quickly. To that end it could be wise to start out only with ones lonely operation to then dare moving on to wider applications. The changing connotations of such words belonging to the dom of the second estate speaks of a bourgeoisie revolt against the knights occuring towards the end of the middle ages but nevertheless not finally sealed until the catastrophe of the world wars, thereby destroying the remnants of legitimate authority in any capacity of earnest ealdormen doomhood, the consigned consigliere henceforth residing but abiding in the blazing bowels of mount doom. Meseems the neat manyfold seemliness of ere lands, fielded with unseamed gilden fleece, beflees the areless naysaying weasely fleas. 

Your last point about the monarch reminds me of an amusing anecdote about a former king of Sweden writing an admonishing letter to his youngest son signed off with ``Head of Family, Father, Friend and King" There is no arguing with that, it is just a fact.


\hfill

\texttt{Cologero on 2017-08-06 at 19:41 said: }

Well, this explains why the grammarian was so important to the Anarch!


\hfill

\texttt{Sparrow on 2017-08-31 at 21:55 said: }

It was rather hard to read Stirner, and the obnoxious behavior of some of his followers (as well as their criminal tendencies) further made it difficult to want to read such a figure. But it seems that Junger and Evola managed to find something profound in him. 

I know it is probably impudent to ask, but would you consider writing on the Japanese Kokugaku movement, due to figures such as Motoori Norinaga and Hirata Atsutane attempting to seek a pure Japanese Tradition through the Kojiki and mystical experiences.


\hfill

\texttt{Fax Fay Faz on 2017-09-19 at 18:14 said: }

I believe you're a bit off on Venator's name. Juenger was a close friend of Heidegger and modeled his narrator more on Heidegger than on himself. The character of Attila (a Hun aka ww1 German soldier) seems rather to be the author's self-portrait, or self-caricature.

Venator is born as Martin (Heidegger) but is renamed Manuel (Kant) by his employer. To paraphrase Nietzsche, Kantianism is the philosophy of the civil servant. The Condor (probably Pinochet) wants to ensure Venator's reliability by making him heir to a ``safer" school if thought.

Much worthwhile in this blog, it appears. Please continue with it, and of course leave the archives up.


\end{sffamily}\end{footnotesize}
